\documentclass[11pt]{article}

\usepackage[a4paper,margin=2.2cm]{geometry}
\usepackage[T1]{fontenc}
\usepackage{lmodern}
\usepackage{microtype}
\usepackage{amsmath,amssymb}
\usepackage{booktabs,longtable,array}
\usepackage{enumitem}
\usepackage[colorlinks,linkcolor=black,urlcolor=blue!50!black]{hyperref}

\setlist[itemize]{itemsep=1pt,topsep=2pt,leftmargin=*}
\newcommand{\E}{\mathbb{E}}
\newcommand{\indep}{\perp\!\!\!\perp}

\title{Presenter Guide for the Slides on Neyman Orthogonality}
\author{}
\date{October 2026}

\begin{document}
\maketitle
\tableofcontents

\section{Symbols and how to read them}

\subsection{Conventions}

\begin{itemize}
\item A subscript $0$ marks the true value: $\theta_0$ is read ``theta zero'' or ``theta naught'', $g_0$ ``g zero''.
\item A hat marks an estimate: $\hat\theta$ is read ``theta hat'', $\hat\eta_k$ ``eta hat k'', the nuisance estimated without fold $k$.
\item A subscript $\xi$ marks a quantity of a simulated data-generating process in pretraining: $g_\xi$ is ``g of xi'', the outcome regression of the simulated process $\xi$.
\item A subscript $\Pi$ marks a posterior mean under the prior $\Pi$: $\hat g_\Pi$ is ``g hat pi'', the posterior mean of the outcome regression.
\item Letters in calligraphic font are sets or classes: $\mathcal T_N$ ``script T N'', $\mathcal H$ ``script H'', $\mathcal D_n$ ``script D n'' for the data set.
\end{itemize}

\subsection{Greek letters}

\begin{longtable}{@{}l l p{10.2cm}@{}}
\toprule
Symbol & Name & Meaning on the slides\\
\midrule
\endhead
$\theta$, $\theta_0$, $\hat\theta$ & theta & target parameter, its true value, its estimate; $\theta(g)$ is the target computed at a regression $g$\\
$\eta$, $\eta_0$ & eta & nuisance (all unknown functions besides the target), and its true value\\
$\alpha$, $\alpha_0$ & alpha & Riesz representer; for the ATE $\alpha_0=D/m_0-(1-D)/(1-m_0)$, the inverse propensity weights\\
$\nu$ & nu & the map that defines the target, $\theta(g)=\E[\nu(W;g)]$; for the ATE $\nu(W;g)=g(1,X)-g(0,X)$\\
$\psi$ & psi & score, the function inside the estimating equation $\E[\psi]=0$; $\psi_{\mathrm{DR}}$ is the AIPW score, $\psi_{\mathrm{IV}}$ and $\psi_{\mathrm{PO}}$ the scores of the partially linear model\\
$\varphi$, $\varphi_\xi$ & phi (curly) & efficient influence function (EIF), or a gradient\\
$\varepsilon$ & epsilon & outcome noise in the partially linear model\\
$\rho$ & rho & residual in a conditional moment restriction $\E[\rho\mid Z]=0$\\
$\omega$, $\varrho$ & omega, rho (curly) & weight density of the average derivative, and conditional density of $D$ given $X$\\
$\delta$ & delta & size of a shift of the outcome regression, $Y\mapsto Y+\delta h(Z)$\\
$\delta_N$, $\Delta_N$, $\lambda_N$ & delta, Delta, lambda & sequences tending to zero in the conditions of the DML theorem; $\lambda_N$ bounds the derivative under near-orthogonality\\
$\lambda$ & lambda & ridge penalty in the Gaussian simulation (appendix)\\
$\xi$ & xi & a data-generating process drawn from the pretraining prior\\
$\Pi$ & capital pi & prior over data-generating processes\\
$\kappa$ & kappa & coefficient of the shift along the Riesz representer in the prior correction, $g_\xi+\kappa\alpha_\xi$\\
$\sigma_n^2$, $\sigma^2$ & sigma squared & variance of $\kappa$ in the prior correction; asymptotic variance in the DML theorem\\
$\varpi_i$ & pi (curly, ``varpi'') & Bayesian bootstrap weights in the label correction\\
$\chi$ & chi & parameters of a trained network $F_\chi$\\
$\tau$, $\tau_\xi$ & tau & conditional average treatment effect function $g(1,x)-g(0,x)$; in $F_\tau$ the mixing weight of a path\\
$\gamma$, $\beta$ & gamma, beta & coefficients of the treatment and outcome equations in the Gaussian simulation\\
\bottomrule
\end{longtable}

\subsection{Other symbols}

\begin{longtable}{@{}l p{4cm} p{7.4cm}@{}}
\toprule
Symbol & Read as & Meaning\\
\midrule
\endhead
$W=(Y,Z)$, $Z=(D,X)$ & & observation: outcome $Y$, regressors $Z$, treatment $D$, covariates $X$\\
$g_0(z)=\E[Y\mid Z=z]$ & ``g zero of z'' & outcome regression\\
$m_0$, $\ell_0$, $f_0$ & ``m zero'', ``ell zero'', ``f zero'' & propensity score or treatment regression $\E[D\mid X]$; $\E[Y\mid X]$; nonparametric part of the partially linear model\\
$V$ & & treatment residual $D-m_0(X)$\\
$\E$, $\E_P$ & ``expectation (under $P$)'' & population mean\\
$\E[\,\cdot\mid Z]$ & ``given $Z$'' & conditional expectation\\
$\mathbb P_n$, $\E_{n,k}$ & ``P n'', ``E n k'' & average over the whole sample, average over fold $I_k$\\
$\langle f,g\rangle$ & ``inner product of f and g'' & $\E[f(Z)g(Z)]$\\
$\|f\|$ & ``norm of f'' & $\sqrt{\E[f(Z)^2]}$, the root mean squared size\\
$L^2(P_Z)$ & ``L two of P Z'' & functions of $Z$ with finite mean square\\
$M(\theta,\eta)$ & ``moment map'' & $\E_P[\psi(W;\theta,\eta)]$, the expected score at a fixed nuisance\\
$J_0$ & ``J zero'', the Jacobian & derivative of $M$ in $\theta$ at the truth\\
$\partial_r$, $\partial_\delta$ & ``partial derivative in r'' & derivative in a scalar parameter\\
$\mathrm D_r[\eta-\eta_0]$ & ``D r applied to eta minus eta zero'' & Gateaux derivative of the moment map in the direction $\eta-\eta_0$; square brackets mark the direction\\
$\mathcal T_N$ & ``script T N'' & nuisance realisation set, which contains the nuisance estimates with high probability\\
$T_1,T_2,T_3$ & ``T one'' etc. & the three terms of the error decomposition\\
$O_P(a_n)$, $o_P(a_n)$ & ``big O P'', ``little o P'' & bounded in probability by a multiple of $a_n$; smaller than $a_n$ in probability\\
$\rightsquigarrow N(0,\sigma^2)$ & ``converges in distribution to a normal'' & asymptotic normality\\
$\propto$ & ``proportional to'' & equal up to a constant factor\\
$\otimes$ & ``product'' & product of distributions, the mechanisms combined\\
$\indep$ & ``independent of'' & independence\\
$\sim$ & ``distributed as'' & e.g.\ $\kappa\sim N(0,\sigma_n^2)$\\
$\mapsto$ & ``maps to'' & e.g.\ $\delta\mapsto b_n(\delta;h)$, the bias as a function of $\delta$\\
$:=$ & ``is defined as'' & definition\\
$P_{\delta h}$ & ``P delta h'' & law of the data after the outcomes are shifted by $\delta h(Z)$\\
$b_n(\delta;h)$ & ``b n of delta and h'' & bias of the estimator at sample size $n$ along this shift\\
$w_i=n\,\partial\hat\theta/\partial Y_i$ & ``implicit weight of observation i'' & how much the estimate moves when outcome $i$ moves, times $n$\\
$\mathcal D_n$ & ``script D n'' & data set of $n$ observations, the context of an in-context learner\\
$F_\chi$, $a_\chi$ & ``F chi'', ``a chi'' & trained network, its amortisation error\\
\bottomrule
\end{longtable}

\subsection{How to read the main equations}

\begin{enumerate}[leftmargin=*,itemsep=4pt]
\item $\theta(g)=\langle\alpha_0,g\rangle$ and $\theta(g_0+h)-\theta_0=\langle\alpha_0,h\rangle$.\\
``The target at any regression $g$ is the weighted average of $g$ with weights $\alpha_0$. If the regression is wrong by $h$, the target is wrong by the $\alpha_0$-weighted average of $h$.'' This is why $\alpha_0$ is the derivative of the target with respect to the regression.

\item $\psi_{\mathrm{DR}}(W;\theta,g,\alpha)=\nu(W;g)-\theta+\alpha(Z)\{Y-g(Z)\}$.\\
``Plug-in estimate minus the target, plus the residual weighted by the Riesz representer.'' For the ATE: $g(1,X)-g(0,X)-\theta+\alpha(Z)\{Y-g(Z)\}$, the AIPW score.

\item $\frac1n\sum_k\sum_{i\in I_k}\psi(W_i;\theta,\hat\eta_k)=0$.\\
``Solve for $\theta$ the sample average of the score, where each observation uses nuisances fitted without its own fold.'' This is DML with cross-fitting.

\item $\E_{n,k}\psi(W;\theta_0,\hat\eta_k)=T_1+T_2+T_3$.\\
``The score at the estimated nuisance splits into noise at the true nuisance ($T_1$), the fluctuation of the difference ($T_2$), and the bias from the nuisance error ($T_3$).'' $T_3=M(\theta_0,\hat\eta_k)$ is the expected score at the estimated nuisance.

\item $\mathrm D_r[\eta-\eta_0]=\partial_r\E_P[\psi(W;\theta_0,\eta_0+r(\eta-\eta_0))]$ and $\mathrm D_0[\eta-\eta_0]=0$.\\
``Move the nuisance from the truth a fraction $r$ of the way towards $\eta$ and differentiate the expected score in $r$. Orthogonality says that this slope is zero at $r=0$ for every admissible $\eta$.'' In the figure it is the slope of each curve at the origin.

\item $\mathrm D_0[h]=\langle\alpha_0-A,h\rangle$.\\
``With correction weights $A$, the slope in direction $h$ is the gap between $\alpha_0$ and $A$, averaged against $h$.'' It vanishes for all $h$ only when $A=\alpha_0$.

\item $M(\theta_0,g,\alpha)=-\langle g-g_0,\alpha-\alpha_0\rangle$.\\
``The bias of the AIPW score is minus the product of the two nuisance errors.'' This is the mixed bias property; the bias is second order and vanishes if either error is zero.

\item $M_{\mathrm{PO}}=\langle\ell-\ell_0,m-m_0\rangle-\theta_0\|m-m_0\|^2$.\\
``The partialling-out bias contains a square of the treatment error.'' The score is orthogonal, but not doubly robust when $\theta_0\neq0$.

\item $\E[\alpha(Z)^2-2\nu(W;\alpha)]=\|\alpha-\alpha_0\|^2-\|\alpha_0\|^2$.\\
``The Riesz loss equals the squared distance to the true representer, up to a constant.'' Minimising it learns $\alpha_0$ without a formula for it.

\item $\sqrt N(\hat\theta-\theta_0)=-J_0^{-1}N^{-1/2}\sum_i\psi(W_i;\theta_0,\eta_0)+o_P(1)\rightsquigarrow N(0,\sigma^2)$ (appendix).\\
``The scaled error behaves like an average of the score at the true nuisance, so it is asymptotically normal.''

\item $\theta_\xi+\sum_i\varpi_i\varphi_\xi(W_i)$.\\
``The training label of the label correction: the true target of the simulated process plus a randomly weighted average of its efficient influence function on the simulated data.''

\item $g_\xi+\kappa\alpha_\xi$ with $\kappa\sim N(0,\sigma_n^2)$.\\
``The prior correction shifts each simulated outcome regression along its own Riesz representer by a random amount.''

\item $P_{\delta h}$, $b_n(\delta;h)=\E_{\delta h}[\hat\theta]-\theta(P_{\delta h})$, $\partial_\delta b_n(0;h)=0$.\\
``Shift all outcomes by $\delta h(Z)$; the bias of the estimator after the shift is $b_n$. First-order insensitivity says the bias does not change to first order in $\delta$.''

\item $\Pi(\mathrm d\xi_X,\mathrm d\xi_Y\mid\mathcal D_n)\propto\Pi_{X,Y}(\mathrm d\xi_X,\mathrm d\xi_Y)\prod_ip(X_i;\xi_X)\,p(Y_i\mid D_i,X_i;\xi_Y)$.\\
``Under independent mechanisms, the posterior of the covariate and outcome mechanisms is the prior times the covariate and outcome likelihood. The treatment likelihood does not appear.''

\item $\partial_\delta b_n(0;h)=\E_P[n^{-1}\sum_iw_ih(Z_i)]-\langle\alpha_0,h\rangle$ (appendix).\\
``The sensitivity of an estimator equals the expected average of its implicit weights times $h$, minus the true change of the target.'' The estimator is insensitive when its implicit weights act like the Riesz representer.
\end{enumerate}

\section{Speaker notes}

\subsection*{Title slide}
\begin{itemize}
\item Three parts: what Neyman orthogonality is, how the literature generalises it, and a project on estimators without a score.
\item Notation follows the draft of the planned paper, close to Chernozhukov et al.\ (2018) and Chernozhukov, Newey and Singh (2022).
\end{itemize}

\subsection*{Data, target and Riesz representer}
\begin{itemize}
\item All targets are linear functionals of the outcome regression $g_0$; the ATE is the running example.
\item The Riesz representer $\alpha_0$ is the derivative of the target in the regression; for the ATE it is the inverse propensity weight.
\item Errors in $g$ move the target by $\langle\alpha_0,h\rangle$, the source of plug-in bias.
\end{itemize}

\subsection*{Scores, moment conditions and cross-fitting}
\begin{itemize}
\item A score turns the target into an estimating equation $\E[\psi]=0$; three examples, AIPW and the two partially linear scores.
\item The choice of instrument in a conditional moment restriction decides sensitivity and efficiency.
\item DML solves the cross-fitted sample equation, so each observation uses nuisances fitted without it.
\end{itemize}

\subsection*{Three-term decomposition}
\begin{itemize}
\item $T_1$ is the oracle noise, $T_2$ the overfitting term, $T_3$ the nuisance bias.
\item Each term has its own remedy: CLT, cross-fitting, orthogonality and rates.
\item For the plug-in score $T_3=\langle\alpha_0,\hat g-g_0\rangle$ exactly, which is too large for machine learning rates.
\end{itemize}

\subsection*{Definition}
\begin{itemize}
\item Orthogonality: the Gateaux derivative of the expected score in the nuisance vanishes at the truth.
\item It is a joint property of score, law and the set $\mathcal T_N$.
\item Near-orthogonality allows a derivative of order $o(N^{-1/2})$.
\end{itemize}

\subsection*{Moment maps along a nuisance direction}
\begin{itemize}
\item Each curve is the expected score as the nuisance moves away from the truth.
\item Orthogonal scores have zero slope at $r=0$: AIPW in (a), both partially linear scores in (b).
\item Plug-in and IPW in (a) and the instrument $D$ in (b) have a nonzero slope.
\end{itemize}

\subsection*{Construction of orthogonal scores}
\begin{itemize}
\item For linear functionals the only orthogonal correction is the Riesz representer times the residual.
\item In the partially linear model any instrument with $\E[A\mid X]=0$ works, so orthogonal scores are not unique.
\item Historically, Neyman's residual score projects off the nuisance scores.
\end{itemize}

\subsection*{Double robustness and the mixed bias property}
\begin{itemize}
\item Hierarchy: mixed bias implies double robustness, which implies orthogonality.
\item The converse fails; $\psi_{\mathrm{PO}}$ is orthogonal but not doubly robust for $\theta_0\neq0$.
\item For AIPW the bias is minus the product of the two errors, at every sample size.
\end{itemize}

\subsection*{Properties and limits of orthogonal scores}
\begin{itemize}
\item Gains: any learner with $n^{-1/4}$ rates or a product rate, no Donsker conditions, valid Wald intervals, efficiency.
\item Limits: only first-order bias, overlap, identification assumptions, and the definition needs an estimating equation.
\item The last limit is the bridge to Part 3.
\end{itemize}

\subsection*{Approaches by term of the decomposition}
\begin{itemize}
\item The literature can be sorted by which term of the decomposition a method handles.
\item Most methods address the first-order part of $T_3$; cross-fitting and stability address $T_2$.
\end{itemize}

\subsection*{Automatic debiased machine learning}
\begin{itemize}
\item The Riesz loss learns $\alpha_0$ without a formula; Lasso, neural nets and forests.
\item Extensions to functionals of M-estimands and numerical influence functions.
\end{itemize}

\subsection*{Orthogonality without an explicit correction}
\begin{itemize}
\item Targeting, sieves, balancing weights and priors reach the same first-order cancellation without adding a correction term.
\item Priors: invariance along the least-favourable direction, propensity-dependent priors, posterior corrections, generalised Bayes.
\end{itemize}

\subsection*{Learned representations and in-context nuisance learners}
\begin{itemize}
\item Representation learning and TabPFN are used as nuisance learners inside DML.
\item Wang, Secrest and Tchakoute (2026) add a propensity head to CausalPFN and run AIPW and TMLE on it.
\end{itemize}

\subsection*{Causal prior-data fitted networks}
\begin{itemize}
\item Causal PFNs return the effect directly and have no orthogonal score.
\item Evidence of bias and miscoverage, prior uncoverage, and prior-induced confounding.
\item Zhang (2026) pretrains on one-step labels.
\end{itemize}

\subsection*{Debiased prior-data fitted networks}
\begin{itemize}
\item Closest work, Melnychuk, Javurek and Feuerriegel (2026): two corrections inside pretraining.
\item Label correction adds a bootstrap average of the EIF; prior correction shifts the simulated regression along the representer.
\item Their guarantees concern the idealised network; amortisation error is set aside.
\end{itemize}

\subsection*{Estimators without a score}
\begin{itemize}
\item Posterior means, regularised plug-ins, end-to-end networks and in-context learners have no estimating equation.
\item The best network under squared loss is the posterior mean, so the prior and the training target decide its properties.
\item The amortisation error of a trained network need not vanish with $n$.
\end{itemize}

\subsection*{First-order insensitivity}
\begin{itemize}
\item Proposed replacement for orthogonality: shift the outcome regression and ask whether the bias moves to first order.
\item It is the derivative part of local asymptotic unbiasedness and is equivalent to regularity along these shifts.
\item It applies to trained networks, not only to idealised ones.
\end{itemize}

\subsection*{Causal PFNs under independent mechanisms}
\begin{itemize}
\item If the prior draws the treatment mechanism independently, the posterior ignores the treatment likelihood (Bayesian ignorability).
\item The posterior mean is then a Bayesian outcome regression with an uncorrected first-order term.
\item The treatment mechanism enters only through the training target or a coupled prior; both corrections of Melnychuk et al.\ are instances.
\end{itemize}

\subsection*{Project question and motivation}
\begin{itemize}
\item Main question: when is a trained estimator without a score first-order insensitive, and how can we check it from data and network alone?
\item Four sub-questions: the criterion, the diagnostic, Bayes rules, trained networks with corrections.
\end{itemize}

\subsection*{Core contribution and positioning}
\begin{itemize}
\item Four contributions, each with its closest work and what is new.
\item The diagnostic, the imbalance of the implicit weights, needs one backward pass.
\item Positioning against Melnychuk et al.\ (2026): trained networks, not the idealised posterior.
\item If asked for evidence: appendix slides on the simulation and the work plan.
\end{itemize}

\subsection*{Appendix}
\begin{itemize}
\item DML theorem: conditions and asymptotic normality, if asked about rates.
\item Efficient influence functions and formulations: equivalent definitions of orthogonality.
\item Local robustness, higher-order terms, forms of robustness: further literature.
\item Implicit weights and proposed results: the technical content behind the contributions.
\item Simulation: posterior mean biased by 4.8 sd at the aligned truth; the imbalance orders the five estimators by their bias.
\item Work plan and risks, and further research questions.
\end{itemize}

\end{document}
